% Part: second-order-logic
% Chapter: sol-and-sets
% Section: introduction

\documentclass[../../../include/open-logic-section]{subfiles}

\begin{document}

\olfileid{sol}{set}{int}

\olsection{引言}

由于二阶逻辑既能对论域的子集进行量化，也能对函数进行量化，因此可以预期，至少有一部分集合论能够在二阶逻辑中开展。所谓“开展”，是指在二阶逻辑中可以表达集合论的性质与命题，且无需使用任何专门关于集合的非逻辑词汇（例如集合论中的隶属谓词符号）。例如，我们可以定义集合的并、交及子集关系，还可以比较集合的大小，并陈述诸如 Cantor（康托尔）定理这样的结果。

\end{document}